\documentclass{article}
\usepackage[utf8]{inputenc}
\usepackage[T1]{fontenc}
\usepackage{geometry}
\usepackage{graphicx}
\usepackage{float}
\usepackage{hyperref}
\usepackage[english,indonesian]{babel}

\title{Optimasi Pemrosesan Template LaTeX Menggunakan Pendekatan Berbasis AI}
\author{Academic Support Agent}
\date{\today}

\begin{document}
\maketitle

\begin{abstract}

\end{abstract}

\chapter{PENDAHULUAN}


\chapter{TINJAUAN PUSTAKA}


% Original content (commented out):
% % Antes de começar, veja o arquivo LEIAME.txt

\documentclass[10pt,a4paper,twocolumn,english,brazil]{article}
\usepackage{sbsr}

%%%% Este arquivo apresenta as configurações globais para o documento em LaTeX %%%%
%%% abaixo devem ser declaradas as informações do trabalho bem como seus autores e afiliações %%%%

\title{título -- modelo sbsr \LaTeX}

\author{Nome e sobrenomes do Primeiro Autor$^1$, Nome e sobrenomes do Segundo Autor$^2$, Nome e sobrenomes do Terceiro Autor$^3$ e Nome e sobrenomes do Último Autor$^4$}

\address{$^1$Afiliação, endereço e e-mail do Primeiro Autor, $^2$Afiliação, endereço e e-mail do Segundo Autor, $^3$Afiliação, endereço e e-mail do Terceiro Autor;... e $^4$Afiliação, endereço e e-mail do Último Autor}

%%%% A partir deste ponto se inicia o documento %%%%

\begin{document}


% Produz o Título do trabalho (NÃO ALTERAR) 
\twocolumn[
\begin{@twocolumnfalse}
\maketitle
\end{@twocolumnfalse}
]


%%%% Adiciona os arquivos de cada seção ao corpo do trabalho %%%%
%%%% (ALTERAR APENAS SE AS SEÇÕES FOREM INTRODUZIDAS/REMOVIDAS) %%%%

\input{01_Resumo.tex}

\input{02_Introducao.tex}

\input{03_Material_e_metodos.tex}
	
\input{04_Resultados.tex}
	
\input{05_Discussao.tex}

\input{06_Conclusoes.tex}

%%%%%%%%%%%%%%%%%%%%%%%%%%%%%%% ATENÇÃO %%%%%%%%%%%%%%%%%%%%%%%%%%%%%%%%
%%%% O arquivo considerações extras (consideracoes_extras.tex) não faz parte do modelo e deve ser removido do trabalho final. Para isto, comente a linha abaixo %%%%
\input{consideracoes_extras.tex} 

%%%% Produz as referências bibliográficas (NÃO ALTERAR) 
%%%% (MAIS INFORMAÇÕES EM LEIAME.txt)

\renewcommand{\refname}{REFERÊNCIAS}
{\small
\bibliographystyle{unsrt} 
\bibliography{07_SBSR.bib}}

\end{document} % Fim do documento %%%%(NÃO REMOVER)%%%%


\end{document}